% Source: Hatcher, "Algebraic Topology", Chapter 1 "The Fundamental Group"
% (2002 Cambridge University Press edition), PDF pages 30-105.
% Transcribed page-by-page by vision subagents, then merged and restructured into
% leanblueprint conventions (semantic \label, bracketed titles, \uses dependency edges) below.
% This file is built up incrementally in chunks following the book's own subsection
% boundaries; see label_map.json for the running Hatcher-numbering <-> label map.

\section*{The Idea of the Fundamental Group}

For a space $X$ and basepoint $x_0$, the fundamental group consists of loops at
$x_0$ modulo homotopies through loops at $x_0$. Loop concatenation models the
group operation. The complement of one circle has fundamental group $\Z$; the
complement of two unlinked circles has fundamental group $\Z*\Z$; the
complement of two linked circles has fundamental group $\Z\times\Z$.

\begin{example}[The Borromean rings and the commutator $aba^{-1}b^{-1}$]
\label{ex:borromean-rings-commutator}
\group{Fundamental Group}
Three Borromean rings have the property that deleting any one ring unlinks
the other two.  Regarding one ring as a loop in the complement of the other
two, its four successive passages represent
$a+b-a-b=aba^{-1}b^{-1}$.  This commutator is nontrivial for the complement
of two unlinked circles but becomes trivial when those circles are linked.
\end{example}

\section{1.1 Basic Constructions}
\label{sec:basic-constructions}

\subsection*{Paths and Homotopy}

\begin{definition}[Path]
\label{def:path}
\lean{HatcherLib.Path}
\leanok
\group{Fundamental Group}
A path in a space $X$ is a continuous map $f:I=[0,1]\to X$.
\end{definition}

\begin{definition}[Homotopy of paths and homotopic paths]
\label{def:homotopy-of-paths}
\lean{HatcherLib.Path.Homotopy, HatcherLib.Path.Homotopic}
\uses{def:homotopy, def:path}
\leanok
\group{Fundamental Group}

A homotopy of paths in $X$ is a family $f_t:I\to X$, $0\leq t\leq1$, with fixed
endpoints $f_t(0)=x_0$ and $f_t(1)=x_1$, such that
$F:I\times I\to X$, $F(s,t)=f_t(s)$, is continuous. We write $f_0\simeq f_1$
when such a homotopy exists.
\end{definition}

\begin{example}[Linear homotopies in $\R^n$]
\label{ex:linear-homotopies-convex}
\lean{HatcherLib.affinePathHomotopy, HatcherLib.affine_paths_homotopic, HatcherLib.convex_paths_homotopic}
\uses{def:homotopy-of-paths}
\leanok
\dcref{ch1:1.1}
\group{Fundamental Group}

If $f_0,f_1:I\to\R^n$ have the same endpoints, then
$f_t(s)=(1-t)f_0(s)+tf_1(s)$ is a homotopy. The same formula stays in any
convex subspace $X\subset\R^n$, so paths in $X$ with prescribed endpoints are
homotopic.
\end{example}

\begin{proposition}[Homotopy of paths is an equivalence relation]
\label{prop:homotopy-paths-equivalence-relation}
\lean{HatcherLib.pathHomotopic_equivalence}
\uses{def:homotopy-of-paths}
\leanok
\dcref{ch1:1.2}
\group{Fundamental Group}

Homotopy is an equivalence relation on paths with fixed endpoints.
\end{proposition}

\begin{proof}
\sketch Reflexivity uses the constant homotopy, symmetry reverses the parameter,
and transitivity concatenates two homotopies in the $t$-coordinate. Continuity
of the concatenated homotopy follows from the pasting lemma.
\end{proof}

\begin{definition}[Composition (product) of paths]
\label{def:path-composition}
\lean{HatcherLib.pathProduct, HatcherLib.pathProduct_homotopic}
\uses{def:path}
\leanok
\group{Fundamental Group}

If $f(1)=g(0)$, define
\[
(f\cdot g)(s)=\begin{cases}f(2s),&0\leq s\leq\frac12,\\
g(2s-1),&\frac12\leq s\leq1.
\end{cases}
\]
\end{definition}

If $f_0\simeq f_1$ and $g_0\simeq g_1$ through compatible endpoint-fixed
homotopies, then $f_0\cdot g_0\simeq f_1\cdot g_1$ by concatenating the
homotopies at each parameter value.

\begin{definition}[Loop, basepoint, and the set $\pi_1(X,x_0)$]
\label{def:loop-and-pi1-set}
\lean{HatcherLib.Loop, HatcherLib.PiOneSet, HatcherLib.PiOne}
\uses{def:path, def:homotopy-of-paths}
\leanok
\group{Fundamental Group}

A loop at $x_0\in X$ is a path $f:I\to X$ with $f(0)=f(1)=x_0$. The set of
homotopy classes of loops at $x_0$ is denoted $\pi_1(X,x_0)$.
\end{definition}

\begin{proposition}[$\pi_1(X,x_0)$ is a group]
\label{prop:pi1-is-a-group}
\lean{HatcherLib.PiOne, HatcherLib.piOneClass_pathProduct}
\uses{def:loop-and-pi1-set, def:path-composition}
\leanok
\dcref{ch1:1.3}
\group{Fundamental Group}

$\pi_1(X,x_0)$ is a group under $[f][g]=[f\cdot g]$.
\end{proposition}

\begin{proof}
\sketch Reparametrization $f\varphi$ by a continuous $\varphi:I\to I$ fixing
the endpoints is homotopic to $f$ through
$f((1-t)\varphi(s)+ts)$. This proves associativity and identity laws, since
the two parenthesizations and products with constant paths are
reparametrizations. For $\bar f(s)=f(1-s)$, shortening $f$ and concatenating
the shortened path with its reverse gives a homotopy to a constant path;
similarly $\bar f\cdot f$ is constant. Thus $[\bar f]$ is the inverse of
$[f]$.
\end{proof}

\begin{example}[$\pi_1$ of a convex set is trivial]
\label{ex:convex-set-trivial-pi1}
\lean{HatcherLib.convex_piOne_subsingleton, HatcherLib.convex_loops_homotopic}
\uses{ex:linear-homotopies-convex, def:loop-and-pi1-set}
\leanok
\dcref{ch1:1.4}
\group{Fundamental Group}

For a convex $X\subset\R^n$ and $x_0\in X$, $\pi_1(X,x_0)=0$, because every
loop is homotopic to the constant loop by the linear homotopy.
\end{example}

\begin{definition}[Change-of-basepoint map]
\label{def:change-of-basepoint-map}
\lean{HatcherLib.changeBasepoint}
\uses{def:path-composition, def:loop-and-pi1-set}
\leanok
\group{Fundamental Group}

For a path $h$ from $x_0$ to $x_1$, define
$\beta_h:\pi_1(X,x_1)\to\pi_1(X,x_0)$ by
$\beta_h[f]=[h\cdot f\cdot\bar h]$, where $\bar h(s)=h(1-s)$.
\end{definition}

\begin{proposition}[The change-of-basepoint map is an isomorphism]
\label{prop:change-of-basepoint-isomorphism}
\lean{HatcherLib.changeBasepoint, HatcherLib.changeBasepoint_bijective}
\uses{def:change-of-basepoint-map}
\dcref{ch1:1.5}
\group{Fundamental Group}

$\beta_h$ is an isomorphism with inverse $\beta_{\bar h}$.
\end{proposition}

\begin{proof}
\sketch Concatenation modulo homotopy gives
$\beta_h([f][g])=\beta_h[f]\beta_h[g]$. The relations
$h\cdot\bar h\simeq c_{x_0}$ and $\bar h\cdot h\simeq c_{x_1}$ show that
$\beta_h\beta_{\bar h}$ and $\beta_{\bar h}\beta_h$ are identities.
\end{proof}

\begin{definition}[Simply-connected space]
\label{def:simply-connected}
\lean{HatcherLib.SimplyConnected}
\uses{def:loop-and-pi1-set}
\leanok
\group{Fundamental Group}

A space is simply-connected if it is path-connected and its fundamental group
is trivial.
\end{definition}

\begin{proposition}[Simply-connected iff unique homotopy class of connecting paths]
\label{prop:simply-connected-iff-unique-path-class}
\lean{HatcherLib.simplyConnected_iff_unique_path_class}
\uses{def:simply-connected, def:homotopy-of-paths}
\leanok
\dcref{ch1:1.6}
\group{Fundamental Group}

A space $X$ is simply-connected iff every two points of $X$ are joined by a
unique homotopy class of paths.
\end{proposition}

\begin{proof}
\sketch If $\pi_1(X)=0$ and $f,g$ join $x_0$ to $x_1$, then
$f\simeq f\cdot\bar g\cdot g\simeq g$. Conversely, uniqueness for paths from
$x_0$ to itself identifies every loop with the constant loop, while existence
of paths gives path-connectedness.
\end{proof}

\subsection*{The Fundamental Group of the Circle}

\begin{theorem}[The fundamental group of the circle]
\label{thm:fundamental-group-of-circle}
\lean{HatcherLib.geometricCirclePiOneMulEquiv, HatcherLib.geometricCircleWindingLoop, HatcherLib.geometricCirclePiOneMulEquiv_windingLoop}
\uses{def:loop-and-pi1-set, def:covering-space}
\leanok
\dcref{ch1:1.7}
\group{Fundamental Group}

$\pi_1(S^1)$ is infinite cyclic, generated by $[\omega]$ for
$\omega(s)=(\cos2\pi s,\sin2\pi s)$ based at $(1,0)$.
\end{theorem}

\begin{definition}[Covering space]
\label{def:covering-space}
\lean{HatcherLib.CoveringMap, HatcherLib.EvenlyCovered}
\leanok
\group{Fundamental Group}
Given $p:\widetilde X\to X$, every $x\in X$ must have an open neighborhood
$U$ such that $p^{-1}(U)$ is a disjoint union of open sets each mapped
homeomorphically onto $U$.
\end{definition}

\begin{proof}[Proof of \cref{thm:fundamental-group-of-circle}]
\sketch For $p:\R\to S^1$, $p(u)=(\cos2\pi u,\sin2\pi u)$, path and homotopy
lifting hold. A loop based at $(1,0)$ lifts from $0$ to a unique integer $n$;
the lifted path is linearly homotopic in $\R$ to $u\mapsto nu$, so the loop is
homotopic to $\omega_n$. A homotopy between $\omega_m$ and $\omega_n$ lifts
with initial point $0$ and has constant endpoint, forcing $m=n$. Lifting is
proved by subdividing compact parameter intervals into pieces over evenly
covered neighborhoods, lifting successively, and using uniqueness to patch.
\end{proof}

\begin{theorem}[Fundamental Theorem of Algebra]
\label{thm:fundamental-theorem-of-algebra}
\lean{HatcherLib.fundamentalTheoremOfAlgebra}
\uses{thm:fundamental-group-of-circle}
\leanok
\dcref{ch1:1.8}
\group{Fundamental Group}

Every nonconstant polynomial over $\C$ has a complex root.
\end{theorem}

\begin{proof}
\sketch If $p(z)=z^n+a_1z^{n-1}+\cdots+a_n$ has no zero, normalize
$p(re^{2\pi is})/p(r)$ to obtain loops in $S^1$. Varying $r$ makes them
homotopic to the null loop at $r=0$. For $r>1+\sum|a_i|$, the homotopy
$z^n+t(a_1z^{n-1}+\cdots+a_n)$ has no zero on $|z|=r$, so the same loop is
homotopic to degree $n$. The circle theorem forces $n=0$, a contradiction.
\end{proof}

The constructive proof by Sperner is cited in \cite{AignerZiegler1999}.

\begin{theorem}[Brouwer fixed point theorem, dimension 2]
\label{thm:brouwer-fixed-point-dim2}
\lean{HatcherLib.euclideanClosedUnitDisk_exists_fixedPoint}
\uses{thm:fundamental-group-of-circle, def:retraction}
\leanok
\dcref{ch1:1.9}
\group{Fundamental Group}

Every continuous $h:D^2\to D^2$ has a fixed point.
\end{theorem}

\begin{proof}
\sketch If $h$ had no fixed point, the ray from $h(x)$ through $x$ would
define a continuous retraction $D^2\to S^1$. A loop in $S^1$ contracts
linearly in $D^2$; composing with the retraction would contract it in $S^1$,
contradicting $\pi_1(S^1)\cong\Z$.
\end{proof}

\begin{theorem}[Borsuk--Ulam theorem, dimension 2]
\label{thm:borsuk-ulam-dim2}
\lean{HatcherLib.borsukUlam_two_of_continuous}
\uses{thm:fundamental-group-of-circle}
\leanok
\dcref{ch1:1.10}
\group{Fundamental Group}

For every continuous $f:S^2\to\R^2$, there is $x\in S^2$ with $f(x)=f(-x)$.
\end{theorem}

\begin{proof}
\sketch Otherwise $g(x)=(f(x)-f(-x))/|f(x)-f(-x)|$ maps $S^2$ to $S^1$
and satisfies $g(-x)=-g(x)$. A lift on the equator changes by $q/2$ after
half a turn for an odd integer $q$, so the equatorial loop has nonzero degree.
It is also the image under $g$ of a loop contracting in $S^2$, a contradiction.
\end{proof}

\begin{corollary}[Covering $S^2$ by three closed sets]
\label{cor:sphere-three-closed-sets-antipodal}
\lean{HatcherLib.sphere_three_closed_sets_antipodal}
\uses{thm:borsuk-ulam-dim2}
\leanok
\dcref{ch1:1.11}
\group{Fundamental Group}

If $S^2=A_1\cup A_2\cup A_3$ with each $A_i$ closed, some $A_i$ contains
an antipodal pair.
\end{corollary}

\begin{proof}
\sketch Apply Borsuk--Ulam to
$x\mapsto(d(x,A_1),d(x,A_2))$. If either coordinate vanishes, the corresponding
closed set contains the pair; otherwise both points lie in $A_3$.
\end{proof}

\begin{proposition}[$\pi_1$ of a product]
\label{prop:pi1-of-product-space}
\lean{HatcherLib.piOneProdMulEquiv}
\uses{def:loop-and-pi1-set}
\leanok
\dcref{ch1:1.12}
\group{Fundamental Group}

For path-connected $X,Y$,
$\pi_1(X\times Y,(x_0,y_0))\cong\pi_1(X,x_0)\times\pi_1(Y,y_0)$.
\end{proposition}

\begin{proof}
\sketch Coordinate projections identify loops and homotopies in $X\times Y$
with pairs in $X$ and $Y$, and this bijection respects concatenation.
\end{proof}

\begin{definition}[Induced homomorphism]
\label{def:induced-homomorphism}
\lean{HatcherLib.inducedPiOne, HatcherLib.inducedPiOneOfEq}
\uses{def:loop-and-pi1-set}
\group{Fundamental Group}

For a basepoint-preserving map $\varphi:(X,x_0)\to(Y,y_0)$, define
$\varphi_*:\pi_1(X,x_0)\to\pi_1(Y,y_0)$ by $\varphi_*[f]=[\varphi f]$.
Then $(\varphi\psi)_*=\varphi_*\psi_*$ and $\one_*=\one$.
\end{definition}

\begin{lemma}[Loops decompose over a path-connected open cover]
\label{lem:loop-decomposition-open-cover}
\lean{HatcherLib.loop_decomposition_of_pathConnectedOpenCover}
\uses{def:loop-and-pi1-set}
\leanok
\dcref{ch1:1.15}
\group{Fundamental Group}

If $X=\bigcup_\alpha A_\alpha$, each $A_\alpha$ is path-connected and open,
contains $x_0$, and every pairwise intersection is path-connected, then every
loop at $x_0$ is homotopic to a product of loops each contained in one
$A_\alpha$.
\end{lemma}

\begin{proof}
\sketch Subdivide the parameter interval so each subpath lies in one cover set.
At each subdivision point choose a path in the adjacent intersection from $x_0$
to that point, and insert that path and its inverse. The resulting factors are
loops in individual cover sets.
\end{proof}

\begin{proposition}[$\pi_1(S^n) = 0$ for $n \ge 2$]
\label{prop:pi1-sphere-trivial-n-geq-2}
\lean{HatcherLib.sphereSimplyConnected_of_two_le, HatcherLib.sphere_piOne_subsingleton_of_two_le}
\uses{def:induced-homomorphism, lem:loop-decomposition-open-cover, ex:convex-set-trivial-pi1}
\leanok
\dcref{ch1:1.14}
\group{Fundamental Group}

$\pi_1(S^n)=0$ if $n\geq2$.
\end{proposition}

\begin{proof}
\sketch Cover $S^n$ by complements of two antipodal points. Each is
homeomorphic to $\R^n$ and their intersection is $S^{n-1}\times\R$, which is
path-connected for $n\geq2$. The lemma decomposes every loop into loops in
contractible cover sets.
\end{proof}

\begin{corollary}[$\R^2$ is not homeomorphic to $\R^n$ for $n \ne 2$]
\label{cor:r2-not-homeomorphic-rn}
\lean{HatcherLib.euclideanTwo_not_homeomorphic_euclidean}
\uses{prop:pi1-of-product-space, prop:pi1-sphere-trivial-n-geq-2, thm:fundamental-group-of-circle}
\leanok
\dcref{ch1:1.16}
\group{Fundamental Group}

$\R^2$ is not homeomorphic to $\R^n$ for $n\ne2$.
\end{corollary}

\begin{proof}
\sketch For $n=1$, deleting a point leaves $\R^2$ connected but $\R$ disconnected.
For $n>2$, $\R^n-\{x\}\simeq S^{n-1}$ has trivial fundamental group,
whereas $\R^2-\{0\}\simeq S^1$ has fundamental group $\Z$.
\end{proof}

\begin{proposition}[Retracts and deformation retracts induce injections/isomorphisms on $\pi_1$]
\label{prop:retract-induces-injective-pi1}
\lean{HatcherLib.retract_inducedPiOne_injective, HatcherLib.deformationRetractPiOneMulEquiv}
\uses{def:induced-homomorphism, def:retraction, def:deformation-retraction}
\leanok
\dcref{ch1:1.17}
\group{Fundamental Group}

If $A$ is a retract of $X$, inclusion induces an injection $\pi_1(A)\to\pi_1(X)$;
if $A$ is a deformation retract, the map is an isomorphism.
\end{proposition}

\begin{proof}
\sketch For a retraction $r$, $r_*i_*=\one$, so $i_*$ is injective. A deformation
retraction supplies a homotopy from the identity to a map through $A$, proving
surjectivity as well.
\end{proof}

\begin{definition}[Homotopy of maps of pairs; based homotopy equivalence]
\label{def:based-homotopy-equivalence}
\lean{HatcherLib.PairHomotopy, HatcherLib.BasedHomotopy, HatcherLib.basedHomotopyEquivType, HatcherLib.BasedHomotopyEquiv.toHomotopyEquiv}
\uses{def:homotopy}
\leanok
\group{Fundamental Group}

A homotopy of pairs carries the first subspace into the second at every time.
It is based when it fixes the basepoint. A based homotopy equivalence consists
of mutually inverse based maps up to based homotopy.
\end{definition}

\begin{proposition}[Homotopy equivalences induce isomorphisms on $\pi_1$, basepoints unconstrained]
\label{prop:homotopy-equivalence-induces-pi1-iso}
\lean{HatcherLib.homotopyEquivPiOneMulEquiv}
\uses{def:induced-homomorphism, lem:basepoint-change-homotopy, def:homotopy-equivalence}
\leanok
\dcref{ch1:1.18}
\group{Fundamental Group}

If $\varphi:X\to Y$ is a homotopy equivalence, then
$\varphi_*:\pi_1(X,x_0)\to\pi_1(Y,\varphi(x_0))$ is an isomorphism.
\end{proposition}

\begin{lemma}[Basepoint tracking under an unbased homotopy]
\label{lem:basepoint-change-homotopy}
\lean{HatcherLib.inducedPiOne_homotopy_eq_changeBasepoint}
\uses{def:change-of-basepoint-map, def:induced-homomorphism}
\leanok
\dcref{ch1:1.19}
\group{Fundamental Group}

If $\varphi_t:X\to Y$ is a homotopy and $h(t)=\varphi_t(x_0)$, then
$\varphi_{0*}=\beta_h\varphi_{1*}$.
\end{lemma}

\begin{proof}
\sketch For a loop $f$ at $x_0$, the loops
$h_t\cdot(\varphi_t f)\cdot\bar h_t$, with $h_t$ the initial segment of $h$,
form a homotopy based at $\varphi_0(x_0)$.
\end{proof}

\begin{proof}[Proof of \cref{prop:homotopy-equivalence-induces-pi1-iso}]
\sketch Compose the induced maps of a homotopy inverse with the basepoint
change formula. The composites are change-of-basepoint isomorphisms, so each
induced map is injective and the first is also surjective.
\end{proof}

\section{1.2 Van Kampen's Theorem}
\label{sec:van-kampen-theorem}

\subsection*{Free Products of Groups}

\begin{definition}[Free product of groups]
\label{def:free-product-of-groups}
\lean{HatcherLib.FreeProduct, HatcherLib.FreeProductWord, HatcherLib.freeProductInclusion}
\leanok
\group{Fundamental Group}
A free product $*_\alpha G_\alpha$ consists of finite reduced words
$g_1\cdots g_m$, with nonidentity adjacent letters from distinct factors.
The empty word is the identity, and multiplication is concatenation followed
by reduction.
\end{definition}

\begin{definition}[Free group, basis, rank]
\label{def:free-group-basis-rank}
\lean{HatcherLib.FreeGroupOn, HatcherLib.FreeBasis, HatcherLib.IsFree, HatcherLib.chosenFreeBasis, HatcherLib.freeGroupRank, HatcherLib.freeBasisIndexEquivGenerators, HatcherLib.freeBasis_cardinalMk_eq_rank, HatcherLib.freeGroupAbelianizationBasis}
\uses{def:free-product-of-groups}
\group{Fundamental Group}

A free group is a free product of copies of $\Z$. The corresponding generators
form a basis and their cardinality is the rank. Its abelianization is free
abelian on the same basis.
\end{definition}

\begin{proposition}[Universal property of the free product]
\label{prop:free-product-universal-property}
\lean{HatcherLib.freeProductLift, HatcherLib.freeProductLift_comp_inclusion, HatcherLib.freeProductLift_unique}
\uses{def:free-product-of-groups}
\leanok
\group{Fundamental Group}

Homomorphisms $\varphi_\alpha:G_\alpha\to H$ extend uniquely to
$\varphi:*_\alpha G_\alpha\to H$ by
$\varphi(g_1\cdots g_m)=\prod_i\varphi_{\alpha_i}(g_i)$.
\end{proposition}

\begin{theorem}[Van Kampen's theorem]
\label{thm:van-kampen-theorem}
\uses{prop:free-product-universal-property, lem:loop-decomposition-open-cover, lem:van-kampen-rectangular-subdivision, lem:van-kampen-grid-cell-sweep}
\dcref{ch1:1.20}
\group{Fundamental Group}

Let $X=\bigcup_\alpha A_\alpha$, where the $A_\alpha$ are path-connected open
sets containing $x_0$ and all pairwise intersections are path-connected. The
map $\Phi:*_\alpha\pi_1(A_\alpha)\to\pi_1(X)$ is surjective. If all triple
intersections are path-connected, its kernel is the normal subgroup generated
by $i_{\alpha\beta}(\omega)i_{\beta\alpha}(\omega)^{-1}$ for
$\omega\in\pi_1(A_\alpha\cap A_\beta)$, and
\[
\pi_1(X)\cong *_\alpha\pi_1(A_\alpha)/N.
\]
\end{theorem}

\begin{example}[Wedge sums]
\label{ex:wedge-sum-fundamental-group}
\uses{thm:van-kampen-theorem, def:wedge-sum}
\dcref{ch1:1.21}
\group{Fundamental Group}

Suppose each path-connected $X_\alpha$ has a basepoint $x_\alpha$ that is a
deformation retract of an open neighborhood. Then van Kampen gives
\[
\pi_1\left(\bigvee_\alpha X_\alpha\right)\cong *_\alpha\pi_1(X_\alpha).
\]
Indeed, the standard cover retracts onto the individual $X_\alpha$, while every
multiple intersection retracts onto the wedge of the basepoint neighborhoods.
In particular, $\pi_1(\bigvee_\alpha S^1)\cong *_\alpha\Z$.
\end{example}

\begin{proof}[Proof of \cref{thm:van-kampen-theorem}]
\sketch The decomposition lemma gives surjectivity. A homotopy between two
factorizations can be subdivided into rectangles mapping into cover sets.
Paths in intersections inserted at grid vertices show that moving across one
rectangle either combines adjacent factors or changes the cover set of a loop
in an intersection, exactly the defining relations of $N$.
\end{proof}

\begin{example}[Linking of circles, via van Kampen]
\label{ex:linking-of-circles-fundamental-groups}
\uses{thm:van-kampen-theorem, prop:pi1-of-product-space, prop:pi1-sphere-trivial-n-geq-2, ex:borromean-rings-commutator}
\dcref{ch1:1.23}
\group{Fundamental Group}

The complement of one circle retracts onto $S^1\vee S^2$, hence has group
$\Z$. The complement of two unlinked circles retracts onto
$S^1\vee S^1\vee S^2\vee S^2$, hence has group $\Z*\Z$. The complement of
two linked circles retracts onto $S^2\vee(S^1\times S^1)$, hence has group
$\Z\times\Z$.
\end{example}

\begin{example}[Torus knot complements]
\label{ex:torus-knot-complement-fundamental-group}
\uses{thm:van-kampen-theorem, def:mapping-cylinder}
\dcref{ch1:1.24}
\group{Fundamental Group}

For positive $m,n$, let $X_{m,n}=X_m\cup_{S^1}X_n$ be the union of the mapping
cylinders of the degree-$m$ and degree-$n$ maps. When $m,n$ are coprime,
$X_{m,n}$ is a deformation retract of the complement of the torus knot
$z\mapsto(z^m,z^n)$. Van Kampen gives
\[
G_{m,n}:=\pi_1(X_{m,n})=\langle a,b\mid a^m=b^n\rangle.
\]
If $m=1$ or $n=1$, then $G_{m,n}\cong\Z$. For $m,n>1$, the element
$a^m=b^n$ generates an infinite cyclic subgroup
$C=\langle a^m\rangle=\langle b^n\rangle$, which is the center of $G_{m,n}$,
and $G_{m,n}/C\cong\Z_m*\Z_n$. The quotient recovers $m,n$ up to order from
the order $mn$ of its abelianization and the largest torsion order. Coprimality
is needed for the knot-complement realization, not for the construction of
$X_{m,n}$; in fact $X_{m,n}$ embeds in $\R^3$ only in the coprime case.
\end{example}

\begin{example}[The shrinking wedge of circles]
\label{ex:shrinking-wedge-circles}
\uses{thm:fundamental-group-of-circle}
\dcref{ch1:1.25}
\group{Fundamental Group}

For the union of circles of radius $1/n$ centered at $(1/n,0)$, retractions
to the individual circles give a surjection
$\pi_1(X)\to\prod_{n\geq1}\Z$. Thus $\pi_1(X)$ is uncountable; finite
collapses give free quotients of every finite rank, so it is nonabelian.
For a fuller description, see \cite{CannonConner2000}. By \cite{Shelah1988},
the fundamental group of a path-connected, locally path-connected compact
metric space is either finitely generated or uncountable.
\end{example}

\subsection*{Applications to Cell Complexes}

\begin{proposition}[Attaching cells and $\pi_1$]
\label{prop:attaching-cells-pi1}
\uses{thm:van-kampen-theorem, prop:pi1-sphere-trivial-n-geq-2}
\dcref{ch1:1.26}
\group{Fundamental Group}

\begin{enumerate}
\item[(a)] Suppose $Y$ is obtained from a path-connected $X$ by attaching
2-cells $e_\alpha^2$ along $\varphi_\alpha:S^1\to X$. Choose paths $y_\alpha$
from $x_0$ to $\varphi_\alpha(s_0)$, and let $N$ be normally generated by
$[y_\alpha\varphi_\alpha\bar y_\alpha]$. The inclusion induces a surjection
$\pi_1(X,x_0)\to\pi_1(Y,x_0)$ with kernel $N$, so
$\pi_1(Y,x_0)\cong\pi_1(X,x_0)/N$.
\item[(b)] If $Y$ is obtained from $X$ by attaching cells of one fixed
dimension $n>2$, then $X\hookrightarrow Y$ induces an isomorphism on
$\pi_1$.
\item[(c)] For every path-connected cell complex $X$, the inclusion
$X^2\hookrightarrow X$ induces
$\pi_1(X^2,x_0)\cong\pi_1(X,x_0)$.
\end{enumerate}
\end{proposition}

\begin{proof}
\sketch Thicken each attached cell by a strip and apply van Kampen to the
complement of one interior point in each cell. The intersection generators are
the conjugated attaching loops, proving (a). For $n>2$ the corresponding
intersections retract onto $S^{n-1}$ and have trivial fundamental group,
proving (b). Induct over skeleta; compactness of loops and nullhomotopies
reduces the infinite-dimensional case to a finite skeleton, proving (c).
\end{proof}

\begin{corollary}[Genus distinguishes closed orientable surfaces]
\label{cor:genus-surfaces-distinct}
\uses{prop:attaching-cells-pi1}
\dcref{ch1:1.27}
\group{Fundamental Group}

Closed orientable surfaces $M_g$ and $M_h$ are not homotopy equivalent when
$g\ne h$.
\end{corollary}

\begin{proof}
\sketch The 2-cell presentation is
$\pi_1(M_g)=\langle a_1,b_1,\ldots,a_g,b_g\mid\prod_i[a_i,b_i]\rangle$,
whose abelianization is $\Z^{2g}$. A homotopy equivalence preserves the
abelianization, so it forces $g=h$.
\end{proof}

\begin{corollary}[Every group is $\pi_1$ of some 2-complex]
\label{cor:every-group-realized-as-pi1}
\uses{prop:attaching-cells-pi1, ex:wedge-sum-fundamental-group}
\dcref{ch1:1.28}
\group{Fundamental Group}

For every group $G$ there is a 2-dimensional cell complex $X_G$ with
$\pi_1(X_G)\cong G$.
\end{corollary}

\begin{proof}
\sketch Choose a presentation $G=\langle g_\alpha\mid r_\beta\rangle.
Start with $\bigvee_\alpha S^1_\alpha$ and attach one 2-cell for each relator.
The attaching-cell proposition gives the desired fundamental group.
\end{proof}

\begin{example}[$X_G$ for $G=\Z_n$]
\label{ex:cyclic-group-complex-xg}
\uses{cor:every-group-realized-as-pi1}
\dcref{ch1:1.29}
\group{Fundamental Group}

For $G=\langle a\mid a^n\rangle=\Z_n$, $X_G$ is $S^1$ with a 2-cell
attached by $z\mapsto z^n$. For $n=2$ it is $\RP^2$; for $n>2$ the $n$
local sheets meeting along the attaching circle prevent it from being a
surface.
\end{example}

\section{1.3 Covering Spaces}
\label{sec:covering-spaces}

\begin{definition}[Sheets of a covering space]
\label{def:covering-space-sheets}
\lean{HatcherLib.IsSheet, HatcherLib.isSheet_iff_connectedComponentIn, HatcherLib.coveringSheetNumber, HatcherLib.coveringSheetNumber_isLocallyConstant, HatcherLib.coveringSheetNumber_eq}
\uses{def:covering-space}
\group{Covering Spaces}

For an evenly covered connected open $U\subset X$, the components of $p^{-1}(U)$
mapped homeomorphically to $U$ are the sheets over $U$. Their number is
$|p^{-1}(x)|$, independent of $x\in U$, and constant when $X$ is connected.
\end{definition}

For $S^1\vee S^1$, coverings can be represented by 2-oriented graphs; the
finite case follows from an Eulerian circuit argument and the infinite case
from \cite{Konig1990}.

\begin{definition}[Lift of a map]
\label{def:lift-of-a-map}
\lean{HatcherLib.IsLift}
\uses{def:covering-space}
\leanok
\group{Covering Spaces}

A lift of $f:Y\to X$ through $p:\widetilde X\to X$ is a map
$\widetilde f:Y\to\widetilde X$ satisfying $p\widetilde f=f$.
\end{definition}

\begin{proposition}[Homotopy lifting property]
\label{prop:covering-homotopy-property}
\lean{HatcherLib.existsUnique_liftHomotopy, HatcherLib.exists_path_lift}
\uses{def:covering-space, def:lift-of-a-map, thm:fundamental-group-of-circle}
\leanok
\dcref{ch1:1.30}
\group{Covering Spaces}

For a covering $p:\widetilde X\to X$, a homotopy $f_t:Y\to X$, and a lift
$\widetilde f_0$ of $f_0$, there is a unique lifted homotopy
$\widetilde f_t:Y\to\widetilde X$. In particular, paths and homotopies of
paths lift uniquely after an initial lift is chosen.
\end{proposition}

\begin{proposition}[Covering maps are injective on $\pi_1$]
\label{prop:covering-map-pi1-injective}
\lean{HatcherLib.coveringPathClassMap_injective, HatcherLib.coveringPiOneMap_injective, HatcherLib.mem_coveringImageSubgroup_iff_monodromy_fixed}
\uses{prop:covering-homotopy-property, def:loop-and-pi1-set}
\leanok
\dcref{ch1:1.31}
\group{Covering Spaces}

For a pointed covering $p:(\widetilde X,\widetilde x_0)\to(X,x_0)$, $p_*$
is injective. Its image consists exactly of loop classes whose lifts from
$\widetilde x_0$ are loops.
\end{proposition}

\begin{proof}
\sketch Lift a nullhomotopy of the projection of a loop; its lift joins the
original loop to a constant loop, proving injectivity. A loop with a closed
lift is visibly in the image, and homotopy lifting gives the converse.
\end{proof}

\begin{proposition}[Sheets equal the index of the image subgroup]
\label{prop:covering-sheets-equals-index}
\lean{HatcherLib.coveringFiberCosetEquiv, HatcherLib.coveringFiber_cardinal_eq_coset_cardinal, HatcherLib.coveringSheets_eq_subgroup_index}
\uses{def:covering-space-sheets, prop:covering-map-pi1-injective}
\leanok
\dcref{ch1:1.32}
\group{Covering Spaces}

For a pointed covering of path-connected spaces, the number of sheets equals
the index of $H=p_*(\pi_1(\widetilde X,\widetilde x_0))$ in $\pi_1(X,x_0)$.
\end{proposition}

\begin{proof}
\sketch Map a left coset $H[g]$ to the endpoint of the lift of $g$ from
$\widetilde x_0$. The lifting criterion for loops gives well-definedness and
injectivity; path-connectedness of $\widetilde X$ gives surjectivity.
\end{proof}

\begin{proposition}[Lifting criterion]
\label{prop:covering-lifting-criterion}
\lean{HatcherLib.exists_lift_iff_piOne_range_le}
\uses{prop:covering-homotopy-property, def:lift-of-a-map}
\leanok
\dcref{ch1:1.33}
\group{Covering Spaces}

Let $Y$ be path-connected and locally path-connected. A pointed map
$f:(Y,y_0)\to(X,x_0)$ lifts through
$p:(\widetilde X,\widetilde x_0)\to(X,x_0)$ iff
$f_*(\pi_1(Y,y_0))\subset p_*(\pi_1(\widetilde X,\widetilde x_0))$.
\end{proposition}

\begin{proof}
\sketch Necessity follows from $f_*=p_*\widetilde f_*$. For sufficiency, send
$y$ to the endpoint of the lift of $f\gamma$ for a path $\gamma$ from $y_0$
to $y$. The subgroup condition makes this independent of $\gamma$; local
path-connected neighborhoods inside evenly covered sets prove continuity.
\end{proof}

Local path-connectedness of $Y$ is essential; without it the criterion fails,
as shown by the counterexample in the corresponding exercise.

\begin{proposition}[Unique lifting property]
\label{prop:covering-unique-lifting}
\lean{HatcherLib.lifts_eq_of_eq_at}
\uses{def:lift-of-a-map}
\leanok
\dcref{ch1:1.34}
\group{Covering Spaces}

If $Y$ is connected and two lifts of $f:Y\to X$ agree at one point, they agree
everywhere.
\end{proposition}

\begin{proof}
\sketch The equalizer of the two lifts is open and closed: near an agreement
point both images lie in one sheet, while distinct sheets separate a point of
disagreement.
\end{proof}

\subsection*{The Classification of Covering Spaces}

\begin{definition}[Semilocally simply-connected]
\label{def:semilocally-simply-connected}
\lean{HatcherLib.SemilocallySimplyConnected}
\leanok
\group{Covering Spaces}
$X$ is semilocally simply-connected if every $x\in X$ has a neighborhood $U$
whose inclusion induces the trivial map $\pi_1(U,x)\to\pi_1(X,x)$.
\end{definition}

\begin{proposition}[Construction of the universal cover]
\label{prop:universal-cover-construction}
\lean{HatcherLib.UniversalCoverConstruction.universalCoverEndpoint_isUniversalCoveringMap}
\uses{def:semilocally-simply-connected, prop:simply-connected-iff-unique-path-class}
\leanok
\group{Covering Spaces}

If $X$ is path-connected, locally path-connected, and semilocally
simply-connected, then $X$ has a simply-connected covering space.
\end{proposition}

\begin{proof}
\sketch Let $\widetilde X$ be homotopy classes of paths starting at $x_0$,
with endpoint projection. For a path-connected open $U$ with trivial image of
$\pi_1(U)$, the sets of classes $[\gamma\eta]$, with $\eta$ a path in $U$,
form evenly covered basic neighborhoods. Truncating paths makes
$\widetilde X$ path-connected. A loop at the constant path lifts to a loop
only when it is nullhomotopic, so covering injectivity gives
$\pi_1(\widetilde X)=0$.
\end{proof}

\begin{example}[Universal cover of $X_{m,n}$ by gluing]
\label{ex:xmn-universal-cover-construction}
\uses{ex:torus-knot-complement-fundamental-group, def:covering-space}
\dcref{ch1:1.35}
\group{Covering Spaces}

For $m,n\geq2$, $X_{m,n}$ is the union of the degree-$m$ and degree-$n$
mapping cylinders; $X_{2,2}$ is the Klein bottle. Their universal covers are
$C_m\times\R$ and $C_n\times\R$, where $C_k$ is a $k$-pronged star. Gluing
copies in the alternating tree pattern gives $T_{m,n}\times\R$; successive
finite subtrees retract to the preceding stage, so this space is contractible
and simply-connected. The evident map from each lifted cylinder to its base
cylinder is therefore the universal covering of $X_{m,n}$.
\end{example}

\begin{proposition}[Covering spaces realizing arbitrary subgroups]
\label{prop:covering-space-existence-for-subgroup}
\lean{HatcherLib.exists_coveringSpace_for_subgroup}
\uses{prop:universal-cover-construction}
\leanok
\dcref{ch1:1.36}
\group{Covering Spaces}

Under the hypotheses of \cref{prop:universal-cover-construction}, every
subgroup $H\subset\pi_1(X,x_0)$ is the image of $\pi_1(X_H,\widetilde x_0)$
for some covering $X_H\to X$.
\end{proposition}

\begin{proof}
\sketch On the universal cover identify path classes with the same endpoint
when their concatenated loop represents an element of $H$. The quotient
inherits evenly covered basic neighborhoods, and a loop lifts closed exactly
when its class lies in $H$.
\end{proof}

\begin{proposition}[Basepointed covering spaces are isomorphic iff same image subgroup]
\label{prop:covering-space-isomorphism-criterion}
\lean{HatcherLib.basedCoveringEquivType, HatcherLib.nonempty_basedCoveringEquiv_iff_piOneRanges_eq}
\uses{prop:covering-lifting-criterion, prop:covering-unique-lifting}
\leanok
\dcref{ch1:1.37}
\group{Covering Spaces}

For path-connected, locally path-connected $X$, pointed connected coverings
$p_i:\widetilde X_i\to X$ are pointed-isomorphic iff
$p_{1*}\pi_1(\widetilde X_1)=p_{2*}\pi_1(\widetilde X_2)$.
\end{proposition}

\begin{proof}
\sketch Necessity follows by composing with the isomorphism. For sufficiency,
the lifting criterion gives pointed maps in both directions over $X$; unique
lifting makes their composites identities.
\end{proof}

\begin{theorem}[Classification of covering spaces]
\label{thm:covering-space-classification}
\lean{HatcherLib.BasedConnectedCover.classificationEquiv, HatcherLib.nonempty_coveringEquiv_iff_conj}
\uses{prop:covering-space-existence-for-subgroup, prop:covering-space-isomorphism-criterion}
\dcref{ch1:1.38}
\group{Covering Spaces}

For path-connected, locally path-connected, semilocally simply-connected $X$,
pointed connected covering spaces correspond bijectively to subgroups of
$\pi_1(X,x_0)$. Unpointed isomorphism classes correspond to conjugacy classes
of subgroups.
\end{theorem}

\begin{proof}
\sketch Existence and pointed uniqueness are the preceding propositions.
Changing the basepoint in a fiber conjugates the image subgroup by the
projected connecting loop, and every conjugate is obtained this way.
\end{proof}

\begin{definition}[Universal cover]
\label{def:universal-cover}
\lean{HatcherLib.IsUniversalCoveringMap, HatcherLib.IsUniversalCoveringMap.existsUnique_mapToCover, HatcherLib.IsUniversalCoveringMap.existsUnique_coveringMapToCover, HatcherLib.universalCoverBasedEquivOfIsUniversal}
\uses{thm:covering-space-classification, prop:covering-lifting-criterion}
\group{Covering Spaces}

A simply-connected covering of a path-connected, locally path-connected space
is its universal cover. It covers every connected covering and is unique up
to isomorphism.
\end{definition}

For a covering $p:\widetilde X\to X$, lifting paths gives an action of
$\pi_1(X,x_0)$ on the fiber over $x_0$. Conversely, an action on a discrete set
$F$ produces the quotient of $\widetilde X_0\times F$ by
$([\gamma],\xi)\sim([\lambda\gamma],\rho(\lambda)\xi)$, and every covering
arises in this way. Thus $n$-sheeted coverings correspond to conjugacy classes
of homomorphisms $\pi_1(X,x_0)\to\Sigma_n$.

\begin{definition}[Deck transformation group]
\label{def:deck-transformation-group}
\lean{HatcherLib.DeckTransformationGroup, HatcherLib.deckTransformation_ext, HatcherLib.deckTransformationGroup_isFree}
\uses{def:covering-space, prop:covering-unique-lifting}
\group{Covering Spaces}

Deck transformations of $p:\widetilde X\to X$ are self-isomorphisms of the
covering. They form $G(\widetilde X)$; on connected $\widetilde X$, a deck
transformation is determined by one point and only the identity fixes a point.
\end{definition}

\begin{definition}[Normal covering space]
\label{def:normal-covering-space}
\lean{HatcherLib.IsNormalCovering}
\uses{def:deck-transformation-group}
\group{Covering Spaces}

A connected covering is normal (regular) if any two lifts of each $x\in X$
are exchanged by a deck transformation.
\end{definition}

\begin{proposition}[Deck group and normalizers]
\label{prop:normal-covering-deck-group}
\lean{HatcherLib.isNormalCovering_iff_piOneImageSubgroup_normal, HatcherLib.normalizerPiOneCoverDeckGroupMulEquivOfLocPathConnectedBase, HatcherLib.normalCoverPiOneDeckGroupMulEquiv, HatcherLib.universalCoverPiOneDeckGroupMulEquiv}
\uses{def:normal-covering-space, thm:covering-space-classification, prop:covering-lifting-criterion}
\leanok
\dcref{ch1:1.39}
\group{Covering Spaces}

Let $p:(\widetilde X,\widetilde x_0)\to(X,x_0)$ be a path-connected covering
of a path-connected, locally path-connected space, and set
$H=p_*\pi_1(\widetilde X,\widetilde x_0)$. The covering is normal iff
$H\trianglelefteq\pi_1(X,x_0)$, and
$G(\widetilde X)\cong N(H)/H$. In the normal case this is
$\pi_1(X,x_0)/H$; for the universal cover it is $\pi_1(X,x_0)$.
\end{proposition}

\begin{proof}
\sketch Changing a lift of $x_0$ conjugates $H$. The lifting criterion gives
a deck transformation between two lifts exactly when the corresponding
conjugate subgroup equals $H$, yielding normality and the normalizer. The map
$N(H)\to G(\widetilde X)$ has kernel precisely the loops lifting to loops,
namely $H$.
\end{proof}

\begin{definition}[Group action on a space; covering space action; free action]
\label{def:group-action}
\lean{HatcherLib.GroupAction, HatcherLib.GroupAction.IsCovering, HatcherLib.GroupAction.IsFree, HatcherLib.GroupAction.IsCovering.isFree, HatcherLib.rationalTranslationAction_isFree, HatcherLib.rationalTranslationAction_not_isCovering, HatcherLib.deckTransformationGroup_isCoveringSpaceAction, HatcherLib.subgroup_deckTransformationGroup_isCoveringSpaceAction}
\group{Covering Spaces}
An action $G\to\operatorname{Homeo}(Y)$ is a covering space action if each
$y\in Y$ has a neighborhood $U$ with $g(U)\cap U\neq\emptyset$ only for
$g=\one$. It is free if only $\one$ fixes a point; a covering space action is
free, but a free action need not be a covering space action. For example, an
irrational rotation acts freely on $S^1$ but has dense, rather than discrete,
orbits and is not a covering space action.
\end{definition}

\begin{definition}[Orbit space]
\label{def:orbit-space}
\lean{HatcherLib.OrbitSpace, HatcherLib.orbitProjection, HatcherLib.normalCoverOrbitHomeomorph}
\uses{def:group-action}
\group{Covering Spaces}

The orbit space $Y/G$ has points the orbits $Gy=\{g(y):g\in G\}$. For a
normal covering, $\widetilde X/G(\widetilde X)=X$.
\end{definition}

\begin{proposition}[Covering space actions give normal covering spaces]
\label{prop:covering-space-action-quotient}
\lean{HatcherLib.orbitProjection_isNormalCovering, HatcherLib.actionDeckTransformationMulEquiv, HatcherLib.actionFundamentalGroupQuotientMulEquiv}
\uses{def:group-action, def:orbit-space, prop:normal-covering-deck-group, prop:covering-unique-lifting}
\dcref{ch1:1.40}
\group{Covering Spaces}

If $G$ acts on $Y$ as a covering space action, then $Y\to Y/G$ is normal,
$G(Y)\cong G$ when $Y$ is path-connected, and, when $Y$ is locally
path-connected, $G\cong\pi_1(Y/G)/p_*\pi_1(Y)$.
\end{proposition}

\begin{proof}
\sketch The translates of a neighborhood satisfying the action condition are
disjoint sheets over its image, proving the covering and normality. A deck
transformation agrees with the unique group element carrying one chosen point
to its image. The quotient formula follows from
\cref{prop:normal-covering-deck-group}.
\end{proof}

\begin{example}[The 11-hole torus covers a genus-3 surface]
\label{ex:genus-11-torus-covering}
\uses{prop:covering-space-action-quotient}
\dcref{ch1:1.41}
\group{Covering Spaces}

The 5-fold rotational action on $M_{11}$ is a covering space action with
quotient $M_3$. Thus $\pi_1(M_{11})$ is normal of index $5$ in $\pi_1(M_3)$
with quotient $\Z_5$; the analogous construction gives $M_{mn+1}\to M_{m+1}$.
\end{example}

\begin{example}[Torus and Klein bottle as symmetry-group quotients of the plane]
\label{ex:grid-symmetries-torus-klein-bottle}
\uses{prop:covering-space-action-quotient}
\dcref{ch1:1.42}
\group{Covering Spaces}

Integer translations of $\R^2$ give the torus quotient and symmetry group
$\Z\times\Z$. Adding glide-reflection symmetries of the alternating grid gives
the Klein bottle quotient; its translation subgroup has index two.
\end{example}

\begin{example}[$\pi_1(\RP^n) \cong \Z_2$ for $n\ge 2$]
\label{ex:real-projective-space-covering}
\uses{prop:covering-space-action-quotient, prop:pi1-sphere-trivial-n-geq-2, ex:real-projective-space-cw}
\dcref{ch1:1.43}
\group{Covering Spaces}

The antipodal $\Z_2$-action on simply-connected $S^n$, $n\geq2$, is a
covering space action with quotient $\RP^n$. Therefore
$\pi_1(\RP^n)\cong\Z_2$.
\end{example}

For even-dimensional spheres, only $\Z_2$ can act freely. On odd-dimensional
spheres, cyclic groups act freely by complex rotations, and further noncyclic
rotation groups are classified in \cite{Wolf1984}. Necessary restrictions and
their converse are discussed in \cite{Milnor1957}, \cite{Brown1982},
\cite{MadsenThomasWall1976}, and \cite{DavisMilgram1985}.

\begin{example}[The action of $G_{m,n}$ on the tree $T_{m,n}$]
\label{ex:gmn-action-on-tree}
\uses{ex:xmn-universal-cover-construction, ex:torus-knot-complement-fundamental-group}
\dcref{ch1:1.44}
\group{Covering Spaces}

$G_{m,n}=\langle a,b\mid a^m=b^n\rangle$ acts diagonally on
$T_{m,n}\times\R$. On $\R$, the generators $a,b$ translate by $1/m$ and
$1/n$; this defines a homomorphism $G_{m,n}\to\Z$ after rescaling the cyclic
translation group. On $T_{m,n}$, the kernel is the center
$\langle a^m\rangle=\langle b^n\rangle$, and the quotient $\Z_m*\Z_n$ acts
with vertex stabilizers. If $K$ is the kernel of $G_{m,n}\to\Z$, then $K$
acts freely, hence as a covering space action, on $T_{m,n}$. Therefore
$K\cong\pi_1(T_{m,n}/K)$ is free.
\end{example}

\subsection*{Cayley Complexes}

\begin{definition}[Cayley graph and Cayley complex]
\label{def:cayley-graph-and-complex}
\lean{HatcherLib.CayleyGraph, HatcherLib.cayleyComplexType, HatcherLib.CayleyGenerates, HatcherLib.cayleyGenerates_iff_closure_range_eq_top, HatcherLib.cayleySimpleGraph_connected, HatcherLib.cayleyGraph_left_translate_isLink}
\uses{cor:every-group-realized-as-pi1}
\group{Covering Spaces}

For $G=\langle g_\alpha\mid r_\beta\rangle$, the Cayley graph has vertex set
$G$ and an edge $g\to gg_\alpha$ for each generator. Attach a 2-cell along
each translate of every relator to obtain the Cayley complex $\widetilde X_G$.
Left multiplication by $G$ is a covering space action and
$\widetilde X_G/G=X_G$.
\end{definition}

\begin{example}[Cayley complexes of small groups]
\label{ex:cayley-graph-free-group}
\uses{def:cayley-graph-and-complex}
\dcref{ch1:1.45}
\group{Covering Spaces}

For $G=\Z*\Z$, the Cayley complex is the Cayley graph, the universal cover of
$S^1\vee S^1$.
\end{example}

\begin{example}[Cayley complex of $\Z\times\Z$]
\label{ex:cayley-complex-torus}
\uses{def:cayley-graph-and-complex}
\dcref{ch1:1.46}
\group{Covering Spaces}

For $G=\langle x,y\mid xyx^{-1}y^{-1}\rangle$, the Cayley complex is the
square grid $\R^2$ and the quotient is the torus.
\end{example}

\begin{example}[Cayley complex of $\Z_n$]
\label{ex:cayley-complex-cyclic-group}
\uses{def:cayley-graph-and-complex, ex:real-projective-space-covering}
\dcref{ch1:1.47}
\group{Covering Spaces}

For $G=\langle x\mid x^n\rangle=\Z_n$, the quotient is $S^1$ with a disk
attached by degree $n$, while the Cayley complex is $n$ disks sharing their
boundary circle. For $n=2$ this is $S^2\to\RP^2$.
\end{example}

\begin{example}[Cayley complex of $\Z_2*\Z_2$]
\label{ex:cayley-complex-infinite-dihedral}
\uses{def:cayley-graph-and-complex, def:free-product-of-groups, ex:cayley-complex-cyclic-group}
\dcref{ch1:1.48}
\group{Covering Spaces}

For $G=\Z_2*\Z_2=\langle a,b\mid a^2,b^2\rangle$, the Cayley complex is an
infinite chain of spheres and the quotient is $\RP^2\vee\RP^2$. The same
construction for $\Z_m*\Z_n$ is tree-like.
\end{example}

\section*{Additional Topics}
\subsection{1.A Graphs and Free Groups}
\label{sec:graphs-and-free-groups}

\begin{definition}[Graph, vertices, edges]
\label{def:graph}
\lean{HatcherLib.TopologicalGraph, HatcherLib.TopologicalGraphVertex, HatcherLib.TopologicalGraphEdge, HatcherLib.TopologicalGraph.isEmpty_cell_of_one_lt}
\uses{def:cell-complex}
\group{Fundamental Group}

A graph is a 1-dimensional CW complex with discrete vertex set $X^0$ and
1-cells $e_\alpha$ attached at their endpoints. It has the weak topology with
respect to the cell closures and is locally path-connected.
\end{definition}

\begin{definition}[Subgraph and tree]
\label{def:subgraph-and-tree}
\lean{HatcherLib.TopologicalGraphSubcomplex, HatcherLib.IsTopologicalTree, HatcherLib.IsSpanningTopologicalSubcomplex, HatcherLib.IsMaximalTopologicalTree, HatcherLib.TopologicalGraphSubcomplex.isClosed, HatcherLib.TopologicalGraphSubcomplex.carrier_eq_vertices_union_edges}
\uses{def:graph, def:contractible}
\leanok
\group{Fundamental Group}

A subgraph is a closed union of vertices and edges. A tree is a contractible
graph; a tree in $X$ is maximal when it contains every vertex of $X$.
\end{definition}

\begin{proposition}[Maximal trees exist]
\label{prop:maximal-tree-exists}
\uses{def:subgraph-and-tree}
\dcref{ch1:1A.1}
\group{Fundamental Group}

Every connected graph contains a maximal tree, and every tree extends to one.
\end{proposition}

\begin{proof}
\sketch Starting from a tree, successively attach edges needed to reach every
new vertex, choosing one edge per new vertex. The resulting increasing union
contains all vertices and deformation retracts stage by stage to the original
tree, hence remains a tree.
\end{proof}

\begin{proposition}[$\pi_1$ of a connected graph is free]
\label{prop:fundamental-group-of-graph-is-free}
\uses{prop:maximal-tree-exists, ex:wedge-sum-fundamental-group, prop:quotient-by-contractible-hep-homotopy-equiv}
\dcref{ch1:1A.2}
\group{Fundamental Group}

If $T$ is a maximal tree in a connected graph $X$, then $\pi_1(X)$ is free
with basis one loop for each edge of $X-T$.
\end{proposition}

\begin{proof}
\sketch Collapse $T$ to a point using the quotient homotopy equivalence. The
quotient is a wedge of circles indexed by the edges outside $T$, so the
wedge-sum calculation gives the stated free basis.
\end{proof}

\begin{lemma}[Covering spaces of graphs are graphs]
\label{lem:covering-space-of-graph-is-graph}
\uses{def:graph, prop:covering-homotopy-property}
\dcref{ch1:1A.3}
\group{Covering Spaces}

Every covering space of a graph is a graph, with vertices and edges the lifts
of the vertices and edges of the base graph.
\end{lemma}

\begin{proof}
\sketch Lift each edge interval from every point over its initial endpoint.
Path lifting supplies the cell attachments, and the local homeomorphism
property identifies the resulting CW topology with the covering topology.
\end{proof}

\begin{theorem}[Every subgroup of a free group is free]
\label{thm:subgroup-of-free-group-is-free}
\lean{HatcherLib.subgroup_of_free_is_free}
\uses{prop:covering-space-existence-for-subgroup, prop:covering-map-pi1-injective, lem:covering-space-of-graph-is-graph, prop:fundamental-group-of-graph-is-free}
\dcref{ch1:1A.4}
\group{Covering Spaces}

Every subgroup of a free group is free.
\end{theorem}

\begin{proof}
\sketch Realize the free group as the fundamental group of a wedge of circles.
The subgroup construction gives a covering graph whose fundamental group is
the subgroup by covering injectivity. The graph proposition makes that
fundamental group free.
\end{proof}

\begin{definition}[Edgepath]
\label{def:edgepath}
\uses{def:graph, prop:maximal-tree-exists, lem:covering-space-of-graph-is-graph, def:universal-cover}
\group{Fundamental Group}

An edgepath is a finite concatenation of monotonically traversed edges. In a
tree there is a unique nonbacktracking edgepath between two vertices. Every
homotopy class of paths between vertices in a connected graph contains a
unique nonbacktracking edgepath, obtained by lifting to the tree universal
cover.
\end{definition}

\subsection{1.B $K(G,1)$ Spaces and Graphs of Groups}
\label{sec:kg1-spaces-graphs-of-groups}

\begin{definition}[$K(G,1)$ space]
\label{def:k-g-1-space}
\lean{HatcherLib.IsKGOneSpace}
\uses{def:universal-cover, def:contractible}
\group{Fundamental Group}

A path-connected space with fundamental group $G$ and contractible universal
cover is a $K(G,1)$ space.
\end{definition}

\begin{example}[Graphs, surfaces, and lens spaces as $K(G,1)$'s]
\label{ex:graphs-are-kg1}
\uses{def:k-g-1-space, prop:fundamental-group-of-graph-is-free, lem:covering-space-of-graph-is-graph}
\dcref{ch1:1B.5}
\group{Fundamental Group}

Connected graphs are $K(F,1)$ spaces for free $F$. Also $\RP^\infty$ is a
$K(\Z_2,1)$, $S^\infty/\Z_m$ is a $K(\Z_m,1)$, and products of $K(G,1)$
spaces are $K(G\times H,1)$; in particular $T^n$ is a $K(\Z^n,1)$.
\end{example}

\begin{example}[Knot complements in $S^3$ as $K(G,1)$'s]
\label{ex:knot-complement-in-s3-kg1}
\uses{def:k-g-1-space, ex:xmn-universal-cover-construction}
\dcref{ch1:1B.6}
\group{Fundamental Group}

For a nonempty closed connected $K\subset S^3$, $S^3-K$ is a $K(G,1)$ by
3-manifold theory. For torus knots this follows from the deformation
retraction to $X_{m,n}$ and its contractible universal cover.
\end{example}

\begin{example}[Milnor's construction: $EG$ and $BG$]
\label{ex:milnor-construction-eg-bg}
\uses{def:k-g-1-space, def:group-action, prop:covering-space-action-quotient}
\dcref{ch1:1B.7}
\group{Fundamental Group}

Let $EG$ be the simplicial complex with simplices $[g_0,\ldots,g_n]$ and
faces obtained by deleting entries. It is contractible, and left multiplication
by $G$ is a covering space action. Thus $BG=EG/G$ is a $K(G,1)$ and the
construction is functorial in $G$.
\end{example}

\begin{theorem}[Homotopy type of a $K(G,1)$ is unique]
\label{thm:kg1-unique-homotopy-type}
\uses{def:k-g-1-space, prop:homomorphism-to-kg1-realized-by-map}
\dcref{ch1:1B.8}
\group{Fundamental Group}

The homotopy type of a CW complex $K(G,1)$ is determined by $G$.
\end{theorem}

\begin{proof}
\sketch An isomorphism of fundamental groups is realized in both directions by
the next proposition. The composites induce identity homomorphisms and are
therefore homotopic to the identity by its uniqueness statement.
\end{proof}

\begin{proposition}[Maps into a $K(G,1)$ realize any homomorphism]
\label{prop:homomorphism-to-kg1-realized-by-map}
\uses{def:k-g-1-space, def:cell-complex, prop:maximal-tree-exists}
\dcref{ch1:1B.9}
\group{Fundamental Group}

If $X$ is a connected CW complex and $Y$ is a $K(G,1)$, every homomorphism
$\varphi:\pi_1(X,x_0)\to\pi_1(Y,y_0)$ is induced by a based map $X\to Y$,
unique up to based homotopy.
\end{proposition}

\begin{proof}
\sketch Collapse a maximal tree to the basepoint, realize $\varphi$ on the
1-cells, and extend over 2-cells because their attaching loops map to the
trivial element. Higher attaching maps lift to the contractible universal
cover of $Y$. Apply the same extension argument to $X\times I$ for uniqueness.
\end{proof}

\begin{definition}[Graph of groups; graph product]
\label{def:graph-of-groups}
\lean{HatcherLib.graphOfGroupsType}
\uses{ex:milnor-construction-eg-bg, def:mapping-cylinder}
\leanok
\group{Fundamental Group}

A graph of groups assigns $G_v$ to each vertex of a connected oriented graph
$\Gamma$ and a homomorphism $\varphi_e:G_{\mathrm{tail}(e)}\to
G_{\mathrm{head}(e)}$ to each edge. Glue $K(G_v,1)$ spaces by mapping
cylinders realizing these maps; the fundamental group of the resulting space
$K\Gamma$ is the graph product.
\end{definition}

\begin{theorem}[Graphs of groups with injective edge maps are $K(G,1)$'s]
\label{thm:graph-of-groups-is-kg1}
\uses{def:graph-of-groups, def:k-g-1-space, prop:covering-space-action-quotient, cor:hep-inclusion-homotopy-equiv-deformation-retract}
\dcref{ch1:1B.11}
\group{Fundamental Group}

If every edge homomorphism is injective, then $K\Gamma$ is a $K(G,1)$ and each
$K(G_v,1)\to K\Gamma$ is injective on $\pi_1$.
\end{theorem}

\begin{proof}
\sketch In the universal cover, each lifted mapping cylinder deformation
retracts onto a contractible lift of its target vertex space. Build the
universal cover as an increasing tree of these pieces; each finite stage
retracts to the preceding stage, so the union is contractible. A loop in a
vertex space that is nullhomotopic globally lifts into one vertex-space lift,
proving injectivity.
\end{proof}

\begin{example}[Free products with amalgamation]
\label{ex:free-product-with-amalgamation}
\uses{thm:graph-of-groups-is-kg1, ex:xmn-universal-cover-construction}
\dcref{ch1:1B.12}
\group{Fundamental Group}

For injective maps $C\to A$ and $C\to B$, the graph of groups
$A\leftarrow C\to B$ has fundamental group $A*_C B$, containing $A$ and
$B$. Taking $A=B=C=\Z$ with degree-$m$ and degree-$n$ maps gives
$\langle a,b\mid a^m=b^n\rangle$.
\end{example}

\begin{example}[HNN extensions]
\label{ex:hnn-extension}
\uses{thm:graph-of-groups-is-kg1}
\dcref{ch1:1B.13}
\group{Fundamental Group}

For monomorphisms $\varphi,\psi:C\to A$, the one-vertex one-edge graph of
groups has group generated by $A$ and $t$ with relations
$t\varphi(c)t^{-1}=\psi(c)$. Equal maps give $A\times\Z$, and an automorphism
gives $A\rtimes\Z$.
\end{example}

\begin{example}[Closed surfaces are $K(G,1)$'s]
\label{ex:closed-surfaces-are-kg1}
\uses{thm:graph-of-groups-is-kg1, ex:graphs-are-kg1}
\dcref{ch1:1B.14}
\group{Fundamental Group}

Closed orientable surfaces of genus at least $2$ are graphs-of-groups spaces
with injective edge maps, hence are $K(G,1)$ spaces with universal cover
$\R^2$. The same holds for closed nonorientable surfaces other than $\RP^2$;
nonclosed surfaces retract to graphs.
\end{example}

The correspondence between graphs of groups and group actions on trees is
developed further in \cite{ScottWall1979}.

\begin{example}[A cube graph has free fundamental group on five generators]
\label{ex:cube-graph-fundamental-group-free}
\dcref{ch1:1.22}
\group{Fundamental Group}
\uses{thm:van-kampen-theorem}
Let $X$ be the graph of the twelve edges of a cube, and $T \subset X$ a maximal
tree (a contractible spanning subgraph) using seven of the edges. For each of the
five edges $e_\alpha$ of $X - T$, choose an open neighborhood $A_\alpha$ of $T\cup e_\alpha$ that
deformation retracts onto it; any intersection of two or more $A_\alpha$'s deformation
retracts onto $T$, hence is contractible, and each $A_\alpha \simeq S^1$. By
\cref{thm:van-kampen-theorem}, $\pi_1(X) \cong *_\alpha \pi_1(A_\alpha)$ is free on five generators, one
per edge outside $T$: for each such $e_\alpha$, a loop that goes through $T$ to one end of
$e_\alpha$, across $e_\alpha$, and back through $T$.


\end{example}
\begin{proof}[Proof of \cref{thm:van-kampen-theorem}]
Surjectivity of $\Phi$ is exactly \cref{lem:loop-decomposition-open-cover}. For the
harder half, call a \textbf{factorization} of $[f]\in\pi_1(X)$ a word $[f_1]\cdots[f_k]$ where
each $f_i$ is a loop in some $A_\alpha$ at $x_0$ and $f \simeq f_1\cdots f_k$ in $X$; surjectivity of
$\Phi$ says every $[f]$ has a factorization. Call two factorizations of $[f]$
\textbf{equivalent} if related by a sequence of: (i) combining adjacent terms in the same
$\pi_1(A_\alpha)$ into one term, or (ii) reinterpreting a term $[f_i]\in\pi_1(A_\alpha)$ as lying in
$\pi_1(A_\beta)$ when $f_i$ is a loop in $A_\alpha\cap A_\beta$. Move (i) does not change the
element of $*_\alpha\pi_1(A_\alpha)$ defined by the word; move (ii) does not change its image in
$Q = *_\alpha\pi_1(A_\alpha)/N$. So equivalent factorizations define the same element of $Q$.

It remains to show any two factorizations of $[f]$ are equivalent, which shows
$Q \to \pi_1(X)$ is injective and hence $\ker\Phi = N$ exactly. Given two factorizations
$[f_1]\cdots[f_k]$, $[f_1']\cdots[f_\ell']$, let $F : I\times I \to X$ be a homotopy between the
composed paths. Partition $I\times I$ into rectangles $R_{ij}$ each mapped by $F$ into a
single $A_{ij}$, refining so each point lies in at most three rectangles, and relabel them
$R_1,\ldots,R_{mn}$ in order. For a path $\gamma$ crossing $I\times I$ left to right, $F|\gamma$ is a
loop at $x_0$; let $\gamma_r$ separate the first $r$ rectangles from the rest, so $\gamma_0,\gamma_{mn}$
are the bottom/top edges. At each vertex $v$ of the grid with $F(v)\ne x_0$, choose a
path $g_v$ from $x_0$ to $F(v)$ inside the (path-connected, by hypothesis) intersection
of the $A_{ij}$'s meeting at $v$; inserting $\bar g_vg_v$ at vertices gives a factorization of
$[F|\gamma_r]$ for each $r$, well-defined up to equivalence, and consecutive $\gamma_r,\gamma_{r+1}$
give equivalent factorizations since $F|\gamma_r \simeq F|\gamma_{r+1}$ within a single $A_{ij}$. Choosing
the $g_v$'s along the bottom and top edges compatibly with the given $f_i$'s and
$f_i'$'s shows the factorization of $\gamma_0$ is equivalent to $[f_1]\cdots[f_k]$ and that of
$\gamma_{mn}$ to $[f_1']\cdots[f_\ell']$, so these two are equivalent to each other.
\end{proof}

\begin{example}[The torus]
\label{ex:torus-fundamental-group}
\dcref{ch1:1.13}
\group{Fundamental Group}
\lean{HatcherLib.torusPiOneMulEquiv, HatcherLib.torusWindingLoop, HatcherLib.torusPiOneMulEquiv_windingLoop, HatcherLib.finiteTorusPiOneMulEquiv, HatcherLib.finiteTorusPiOneMulEquiv_windingLoop}
\leanok
\uses{prop:pi1-of-product-space,thm:fundamental-group-of-circle}
By \cref{prop:pi1-of-product-space}, $\pi_1(S^1\times S^1) \approx \Z\times\Z$. A pair
$(p,q) \in \Z\times\Z$ corresponds to a loop $\omega_{p,q}(s) = (\omega_p(s),\omega_q(s))$ winding $p$
times around one $S^1$ factor and $q$ times around the other. This loop can be
knotted, as for $p=3,q=2$ -- the resulting torus knots are studied later
(\cref{ex:torus-knot-complement-fundamental-group}).



More generally, the $n$-dimensional torus (product of $n$ circles) has fundamental
group isomorphic to $\Z^n$, by induction on $n$.
\end{example}
